\documentclass[a4paper,12pt]{article}
\usepackage{cmap}
\usepackage[T2A]{fontenc}
\usepackage[utf8]{inputenc}
\usepackage[russian]{babel}
\usepackage{amsmath,amssymb,mathtools}
\renewcommand{\familydefault}{\sfdefault}
\usepackage{icomma}
\usepackage[a4paper,margin=18mm,top=18mm,bottom=20mm]{geometry}
\usepackage{microtype}
\usepackage{tikz}
\usetikzlibrary{calc,arrows.meta,positioning,shapes.geometric}
\usepackage{enumitem}
\usepackage{booktabs,tabularx,longtable}
\usepackage{xcolor}
\usepackage{setspace}
\usepackage{fancyhdr}
\usepackage[hidelinks]{hyperref}
\usepackage{parskip}

\definecolor{accent}{HTML}{176B87}
\definecolor{darkaccent}{HTML}{0F4C5C}
\definecolor{textgray}{HTML}{2E3438}
\definecolor{softgray}{HTML}{EEF3F5}
\definecolor{warn}{HTML}{A6533B}
\definecolor{ok}{HTML}{386641}

\color{textgray}
\onehalfspacing
\setlength{\parindent}{0pt}
\setlength{\headheight}{15pt}
\setlist[itemize]{leftmargin=6mm,itemsep=2pt,topsep=3pt}
\setlist[enumerate]{leftmargin=7mm,itemsep=5pt,topsep=4pt}

\pagestyle{fancy}
\fancyhf{}
\fancyhead[L]{\small\color{darkaccent}Векторы}
\fancyhead[R]{\small\color{darkaccent}28.09.26}
\fancyfoot[C]{\small\thepage}
\renewcommand{\headrulewidth}{0.3pt}
\renewcommand{\footrulewidth}{0pt}

\newcommand{\vect}[1]{\overrightarrow{#1}}
\newcommand{\key}[1]{\textbf{\color{darkaccent}#1}}
\newcommand{\mistake}[1]{\textbf{\color{warn}#1}}
\newcommand{\good}[1]{\textbf{\color{ok}#1}}

\tikzset{
  vec/.style={-{Stealth[length=3.2mm,width=2.1mm]},line width=1.25pt,color=accent},
  flowstart/.style={ellipse,draw=darkaccent,line width=0.9pt,minimum width=34mm,minimum height=11mm,align=center,fill=softgray},
  flowaction/.style={rounded corners=2mm,draw=darkaccent,line width=0.9pt,text width=48mm,minimum height=13mm,align=center,fill=white},
  flowdecision/.style={diamond,aspect=2.2,draw=darkaccent,line width=0.9pt,text width=39mm,align=center,inner sep=1.5pt,fill=softgray},
  flowarrow/.style={-{Stealth[length=3mm,width=2mm]},line width=1pt,color=darkaccent}
}

\begin{document}
\thispagestyle{empty}
\begin{center}
\vspace*{18mm}
{\Large\color{darkaccent}\textbf{Чек-лист}}\\[6mm]
{\huge\bfseries Векторы}\\[3mm]
{\Large действия и правила построения}\\[14mm]

\begin{tikzpicture}[scale=1.05]
  \coordinate (A) at (0,0);
  \coordinate (B) at (3.2,0.7);
  \coordinate (C) at (5.9,-0.1);
  \draw[vec] (A)--(B) node[midway,above] {$\vec a$};
  \draw[vec] (B)--(C) node[midway,above] {$\vec b$};
  \draw[vec,color=darkaccent,line width=1.6pt] (A)--(C) node[midway,below] {$\vec a+\vec b$};
\end{tikzpicture}

\vfill
{\large Марина Селиверстова}\\[3mm]
{\normalsize 28.09.26}\\[11mm]
{\small\color{darkaccent}Лёвин Артём Александрович эксклюзивно для Марины Селиверстовой}
\end{center}

\begin{tikzpicture}[remember picture,overlay]
\draw[accent!55,line width=1.2pt]
  ([xshift=16mm,yshift=12mm]current page.south west)
  .. controls ([xshift=55mm,yshift=22mm]current page.south west) and ([xshift=-55mm,yshift=4mm]current page.south east) ..
  ([xshift=-16mm,yshift=15mm]current page.south east);
\end{tikzpicture}

\clearpage
\section*{1. Что такое вектор}
\key{Вектор} --- направленный отрезок: у него есть \textbf{начало} и \textbf{конец}.

Если начало --- точка $A$, а конец --- точка $B$, то вектор записывают как $\vect{AB}$. Порядок букв важен: $\vect{AB}$ направлен от $A$ к $B$, а $\vect{BA}$ --- в обратную сторону.

\begin{center}
\begin{tikzpicture}
  \coordinate (A) at (0,0);
  \coordinate (B) at (5.2,0);
  \fill[darkaccent] (A) circle (1.8pt) node[below] {$A$};
  \fill[darkaccent] (B) circle (1.8pt) node[below] {$B$};
  \draw[vec] (A)--(B) node[midway,above] {$\vect{AB}$};
\end{tikzpicture}
\end{center}

\key{Мини-памятка:} первая буква --- начало, вторая --- конец. При переносе вектора в другое место его \textbf{не поворачивают и не растягивают}: длина и направление сохраняются.

\section*{2. Сложение двух векторов: правило треугольника}
Чтобы сложить два вектора, второй вектор переносят параллельно самому себе так, чтобы \key{его начало совпало с концом первого}.

\begin{center}
\begin{tikzpicture}[scale=0.95]
  \coordinate (A) at (0,0);
  \coordinate (B) at (3.2,1.0);
  \coordinate (C) at (6.5,0.2);
  \draw[vec] (A)--(B) node[midway,above] {$\vec a$};
  \draw[vec] (B)--(C) node[midway,above] {$\vec b$};
  \draw[vec,color=darkaccent,line width=1.6pt] (A)--(C) node[midway,below] {$\vec a+\vec b$};
  \fill[darkaccent] (A) circle (1.5pt);
  \fill[darkaccent] (B) circle (1.5pt);
  \fill[darkaccent] (C) circle (1.5pt);
\end{tikzpicture}
\end{center}

\begin{enumerate}
  \item Постройте первый вектор.
  \item К его концу перенесите \textbf{начало} второго вектора, сохраняя его длину и направление.
  \item Соедините начало первого с концом второго. Полученный вектор --- сумма.
\end{enumerate}

\mistake{Типичная ошибка:} совмещать конец первого с концом второго. Нужна связка \textbf{конец первого $\to$ начало второго}.

\section*{3. Вычитание векторов}
Вычитание переводим в сложение:
\[
\vec a-\vec b=\vec a+(-\vec b).
\]

Вектор $-\vec b$ имеет \textbf{ту же длину}, что и $\vec b$, но направлен в противоположную сторону. Иными словами, у него меняются местами начало и конец.

\begin{center}
\begin{tikzpicture}[scale=0.95]
  \coordinate (A) at (0,0);
  \coordinate (B) at (3.2,0.8);
  \coordinate (C) at (5.9,-0.4);
  \draw[vec] (A)--(B) node[midway,above] {$\vec a$};
  \draw[vec] (B)--(C) node[midway,above] {$-\vec b$};
  \draw[vec,color=darkaccent,line width=1.6pt] (A)--(C) node[midway,below] {$\vec a-\vec b$};
\end{tikzpicture}
\end{center}

\key{Главная последовательность:} развернуть вычитаемый вектор $\to$ присоединить его началом к концу первого $\to$ замкнуть путь.

\clearpage
\section*{4. Умножение вектора на число}
Для числа $k$ выполняется
\[
|k\vec a|=|k|\,|\vec a|.
\]
То есть модуль коэффициента задаёт изменение длины, а знак определяет направление:

\begin{itemize}
  \item $k>0$: направление сохраняется;
  \item $k<0$: направление становится противоположным;
  \item $k=0$: получается нулевой вектор $\vec 0$;
  \item дробный коэффициент, например $\frac12$, уменьшает длину вектора в соответствующее число раз.
\end{itemize}

\begin{center}
\begin{tikzpicture}[x=1cm,y=1cm]
  \draw[vec] (0,1.8)--(2.2,1.8) node[midway,above] {$\vec a$};
  \draw[vec] (0,0.8)--(4.4,0.8) node[midway,above] {$2\vec a$};
  \draw[vec] (4.4,-0.2)--(0,-0.2) node[midway,below] {$-2\vec a$};
\end{tikzpicture}
\end{center}

\mistake{Важно:} положительный коэффициент не поворачивает вектор. Отрицательный знак разворачивает направление строго на противоположное.

\section*{5. Сложение нескольких векторов: правило многоугольника}
Для трёх и более векторов действует та же идея: к концу предыдущего вектора присоединяют начало следующего.

\begin{center}
\begin{tikzpicture}[scale=0.9]
  \coordinate (A) at (0,0);
  \coordinate (B) at (2.5,1.1);
  \coordinate (C) at (4.5,0.2);
  \coordinate (D) at (6.8,1.6);
  \draw[vec] (A)--(B) node[midway,above] {$\vec a$};
  \draw[vec] (B)--(C) node[midway,above] {$\vec b$};
  \draw[vec] (C)--(D) node[midway,below] {$\vec c$};
  \draw[vec,color=darkaccent,line width=1.6pt] (A)--(D) node[midway,below] {$\vec a+\vec b+\vec c$};
\end{tikzpicture}
\end{center}

Результирующий вектор соединяет \textbf{начало первого} с \textbf{концом последнего}.

\section*{6. Как выражать один вектор через другие}
Ищите путь от начала нужного вектора к его концу и записывайте последовательность перемещений.

В параллелограмме $ABCD$, если
\[
\vect{AB}=\vec a,\qquad \vect{AD}=\vec b,
\]
то путь $B\to A\to D$ даёт
\[
\vect{BD}=\vect{BA}+\vect{AD}=-\vec a+\vec b=\vec b-\vec a.
\]

\key{Контроль направления:} нужный вектор $\vect{BD}$ должен начинаться в $B$ и заканчиваться в $D$. Это быстрый способ проверить знак и порядок слагаемых.

\section*{7. Ошибки, которые стоит ловить сразу}
\begin{itemize}
  \item \mistake{Перепутаны начало и конец:} $\vect{AB}\neq\vect{BA}$.
  \item \mistake{Второй вектор присоединён концом:} при сложении к концу первого подводят \textbf{начало второго}.
  \item \mistake{При вычитании не развернули второй вектор:} сначала $\vec a-\vec b=\vec a+(-\vec b)$.
  \item \mistake{Вектор повернули произвольно:} при параллельном переносе длина и направление сохраняются.
  \item \mistake{Отрицательный коэффициент изменил только длину:} знак «минус» обязательно меняет направление на противоположное.
  \item \mistake{Результирующий вектор соединяет неверные точки:} соединяйте начало первого и конец последнего.
\end{itemize}

\section*{8. Самопроверка за 20 секунд}
Перед тем как считать построение завершённым, проверьте четыре пункта:
\begin{enumerate}[label=\arabic*.,itemsep=2pt]
  \item При переносе вектор сохранил длину и направление?
  \item При вычитании второй вектор сначала развернут?
  \item Каждый следующий вектор начинается в конце предыдущего?
  \item Результат соединяет начало всего пути с его концом?
\end{enumerate}

\clearpage
\thispagestyle{fancy}
\begin{center}
{\Large\bfseries\color{darkaccent}Алгоритм: сложение и вычитание векторов}
\end{center}
\vspace{6mm}

\begin{center}
\begin{tikzpicture}[node distance=13mm and 17mm]
  \node[flowstart] (start) {Начало};
  \node[flowdecision,below=of start] (op) {Какая операция?};
  \node[flowaction,below left=18mm and 9mm of op] (sum) {Сложение: второй вектор оставьте без изменения};
  \node[flowaction,below right=18mm and 9mm of op] (sub) {Вычитание: замените $\vec a-\vec b$ на $\vec a+(-\vec b)$ и разверните $\vec b$};
  \node[flowaction,below=28mm of op] (attach) {Параллельно перенесите второй вектор: его начало должно совпасть с концом первого};
  \node[flowaction,below=of attach] (close) {Соедините начало первого вектора с концом второго};
  \node[flowstart,below=of close] (finish) {Получен результирующий вектор};

  \draw[flowarrow] (start)--(op);
  \draw[flowarrow] (op)-- node[above left]{\small сложение} (sum);
  \draw[flowarrow] (op)-- node[above right]{\small вычитание} (sub);
  \draw[flowarrow] (sum.south) |- (attach.west);
  \draw[flowarrow] (sub.south) |- (attach.east);
  \draw[flowarrow] (attach)--(close);
  \draw[flowarrow] (close)--(finish);
\end{tikzpicture}
\end{center}

\vfill
\begin{center}
\good{Финальная проверка:} перенос не меняет вектор, а результат начинается там же, где первый вектор, и заканчивается там, где завершился весь путь.
\end{center}

\clearpage
\section*{Тренировочный блок --- Путь А}
\textbf{10 заданий.} Выполняйте построения схематично: важнее сохранить направления, длины в клетках и правильный порядок соединения векторов.

\begin{enumerate}[label=\textbf{\arabic*.}]
  \item Вектор $\vect{AB}$ направлен из точки $A$ в точку $B$. Запишите вектор той же длины, направленный строго противоположно.

  \item Запишите разность $\vec a-\vec b$ в виде суммы двух векторов.

  \item Вектор $\vec p$ имеет длину 4 клетки и направлен вправо. Укажите длину и направление вектора $-2\vec p$.

  \item Вектор $\vec a$ имеет длину 3 клетки и направлен вправо, а $\vec b$ --- 2 клетки вверх. Постройте $\vec a+\vec b$ по правилу треугольника и опишите положение конца результирующего вектора относительно его начала.

  \item Вектор $\vec p$ имеет длину 4 клетки и направлен вправо, а $\vec k$ --- 3 клетки вверх. Постройте $\vec p-\vec k$ и опишите положение его конца относительно начала.

  \item Вектор $\vec p$ имеет длину 2 клетки и направлен вправо, а $\vec k$ --- 2 клетки вверх. Постройте $\vec p+2\vec k$ и опишите положение конца результирующего вектора относительно начала.

  \item Вектор $\vec a$ имеет длину 2 клетки и направлен вправо, а $\vec b$ --- 3 клетки вверх. Постройте $2\vec a+3\vec b$ и опишите итоговое перемещение.

  \item Вектор $\vec p$ имеет длину 3 клетки и направлен вправо, $\vec k$ --- 2 клетки вверх, $\vec t$ --- 1 клетку влево. Постройте $\vec p+\vec k-\vec t$ и опишите итоговое перемещение.

  \item В параллелограмме $ABCD$ задано $\vect{AB}=\vec a$, $\vect{AD}=\vec b$. Выразите $\vect{BD}$ через $\vec a$ и $\vec b$.

  \item В том же параллелограмме выразите $\vect{DB}$ через $\vec a$ и $\vec b$. Проверьте ответ сравнением с $\vect{BD}$.
\end{enumerate}

\clearpage
\section*{Ответы}
\renewcommand{\arraystretch}{1.35}
\begin{longtable}{@{}p{12mm}p{0.82\linewidth}@{}}
\toprule
\textbf{№} & \textbf{Ответ / ориентир для самопроверки}\\
\midrule
\endhead
1 & $\vect{BA}=-\vect{AB}$.\\
2 & $\vec a-\vec b=\vec a+(-\vec b)$.\\
3 & Длина 8 клеток; направление влево.\\
4 & От начала: 3 клетки вправо и 2 клетки вверх. Результат --- $\vec a+\vec b$.\\
5 & Развернуть $\vec k$: $-\vec k$ направлен вниз. От начала: 4 клетки вправо и 3 клетки вниз.\\
6 & $2\vec k$ --- 4 клетки вверх. От начала: 2 клетки вправо и 4 клетки вверх.\\
7 & $2\vec a$ --- 4 клетки вправо, $3\vec b$ --- 9 клеток вверх. Итог: 4 клетки вправо и 9 клеток вверх.\\
8 & $-\vec t$ --- 1 клетка вправо. Итог: 4 клетки вправо и 2 клетки вверх.\\
9 & $\vect{BD}=\vect{BA}+\vect{AD}=-\vec a+\vec b=\vec b-\vec a$.\\
10 & $\vect{DB}=\vec a-\vec b=-\vect{BD}$.\\
\bottomrule
\end{longtable}

\vfill
\begin{center}
{\small\color{darkaccent}Лёвин Артём Александрович эксклюзивно для Марины Селиверстовой}
\end{center}

\end{document}
